\section{Future work}
\label{sec:future}

\paragraph{Completed research runs.}
The immediate next step is to run the pre-registered conditions at research
scale, freeze the dataset and evaluation versions, and replace every placeholder
in this manuscript with reviewed measurements. Negative results will be kept.

\paragraph{Online distillation.}
Sequence-level distillation is a pragmatic baseline. Future work can compare it
with logit matching, intermediate-layer distillation, and on-policy distillation
when the student and teacher share a tokenizer and the teacher fits in memory.
The compute and licensing costs of each must be reported.

\paragraph{Retrieval.}
BM25 is reproducible and dependency-free, but dense retrieval, hybrid retrieval,
reranking, and query rewriting may improve grounding. Future work should measure
whether gains survive on consumer hardware and whether citations remain faithful.
Prompt-injection robustness of retrieved evidence is a priority.

\paragraph{Autocomplete.}
A controlled editor study is needed to measure acceptance, latency sensitivity,
and error modes across Bash, Python, PowerShell, YARA, Sigma, and configuration
files. FIM quality should be compared across base checkpoints with known training
recipes. A language-server integration and a sandboxed evaluation of generated
commands (never automatic execution) are future engineering goals.

\paragraph{Quantization and CPU deployment.}
GGUF conversion support for Qwen3.5 must be validated as llama.cpp evolves.
Future work should compare 4-bit, 5-bit, 8-bit, and mixed-precision deployments
on end-task accuracy, not only perplexity, and should measure laptop-class CPU
latency with realistic context lengths.

\paragraph{Larger controls and scaling.}
The architecture-aware pipeline supports larger Qwen3.5 checkpoints and other
text models. A scaling study could test whether the performance/compute frontier
changes at 1.7B, 4B, and 8B parameters under the same data and evaluation
protocol. This is optional and should not delay the core sub-billion study.

\paragraph{Multimodal security tasks.}
The checkpoint includes a vision tower. Future work could evaluate screenshot
triage, diagram understanding, and phishing-image analysis, but this requires new
data, safety review, and a separate evaluation protocol. It is outside the
current text-only scope.

\paragraph{Human evaluation.}
Task metrics cannot capture whether an explanation is actionable and safe.
Future work should include structured human review with rubrics, inter-annotator
agreement, and disagreement analysis, especially for remediation advice and
detection rules.

\paragraph{Interpretability and forgetting.}
Mechanistic and representation-level analyses could explain why narrow SFT
improves some tasks and degrades others. Measuring forgetting on a broader
retention suite, and testing replay or regularization, would make the efficiency
claim more complete.

\paragraph{Community artifacts.}
The project would benefit from independently reproduced runs, additional
licensed datasets, editor integrations, and adversarial evaluations. All should
follow the repository's provenance and reporting standards.
